% !TeX program = xelatex
\documentclass[UTF8,a4paper,12pt]{ctexart}

\usepackage{amsmath}
\usepackage{amssymb}
\usepackage{geometry}

\geometry{margin=2.5cm}
\setcounter{MaxMatrixCols}{20}

\title{尾巴离地状态下的系统建模}
\author{}
\date{}

\begin{document}

\maketitle

\section{尾巴离地状态下的系统建模}

本文针对轮腿机器人尾巴离地工况建立动力学模型。该工况下，尾巴不与地面接触，因此尾巴不提供额外的地面支撑力，只通过尾电机安装点与机体发生动力学耦合。系统由机体、左右腿、左右轮以及尾巴组成。

\subsection{广义坐标与控制输入}

选取系统广义坐标为
\begin{equation}
    q =
    \begin{bmatrix}
        X_b^h & \phi & \theta_l & \theta_r & \theta_b & \theta_t
    \end{bmatrix}^{\mathrm T},
\end{equation}
其中，$X_b^h$ 为机体水平方向位置，$\phi$ 为偏航角，$\theta_l$ 与 $\theta_r$ 分别为左右腿角度，$\theta_b$ 为机体俯仰角，$\theta_t$ 为尾巴角度。

对应的广义速度和广义加速度分别为
\begin{equation}
    \dot q =
    \begin{bmatrix}
        \dot X_b^h & \dot \phi & \dot \theta_l & \dot \theta_r & \dot \theta_b & \dot \theta_t
    \end{bmatrix}^{\mathrm T},
\end{equation}
\begin{equation}
    \ddot q =
    \begin{bmatrix}
        \ddot X_b^h & \ddot \phi & \ddot \theta_l & \ddot \theta_r & \ddot \theta_b & \ddot \theta_t
    \end{bmatrix}^{\mathrm T}.
\end{equation}

系统控制输入选为
\begin{equation}
    u =
    \begin{bmatrix}
        T_{r \to b} &
        T_{l \to b} &
        T_{wr \to r} &
        T_{wl \to l} &
        T_{t \to b}
    \end{bmatrix}^{\mathrm T},
\end{equation}
其中，$T_{r \to b}$ 和 $T_{l \to b}$ 为左右腿对机体的驱动力矩，$T_{wr \to r}$ 和 $T_{wl \to l}$ 为左右轮对腿部的驱动力矩，$T_{t \to b}$ 为尾电机对机体的力矩。

\subsection{尾巴运动学描述}

设尾电机安装点相对机体质心的前向距离为 $a_{t,p}$，下向距离为 $b_{t,p}$，则安装点相对机体质心的位置可写为
\begin{equation}
    r_{bp}^h = a_{t,p}\cos\theta_b + b_{t,p}\sin\theta_b,
\end{equation}
\begin{equation}
    r_{bp}^v = a_{t,p}\sin\theta_b - b_{t,p}\cos\theta_b.
\end{equation}

尾巴质心相对尾杆方向存在固定偏置角 $\delta_t$，因此定义尾巴质心方向角为
\begin{equation}
    \alpha_t = \theta_t - \theta_b + \delta_t.
\end{equation}

尾电机安装点的水平和竖直加速度分别为
\begin{equation}
    a_{tp}^h =
    a_b^h
    - \left(a_{t,p}\sin\theta_b - b_{t,p}\cos\theta_b\right)\ddot\theta_b
    - \left(a_{t,p}\cos\theta_b + b_{t,p}\sin\theta_b\right)\dot\theta_b^2,
\end{equation}
\begin{equation}
    a_{tp}^v =
    a_b^v
    + \left(a_{t,p}\cos\theta_b + b_{t,p}\sin\theta_b\right)\ddot\theta_b
    - \left(a_{t,p}\sin\theta_b - b_{t,p}\cos\theta_b\right)\dot\theta_b^2.
\end{equation}

尾巴质心加速度为
\begin{equation}
    a_t^h =
    a_{tp}^h
    - l_{t,c}
    \left[
        \left(\ddot\theta_t - \ddot\theta_b\right)\sin\alpha_t
        +
        \left(\dot\theta_t - \dot\theta_b\right)^2\cos\alpha_t
    \right],
\end{equation}
\begin{equation}
    a_t^v =
    a_{tp}^v
    - l_{t,c}
    \left[
        \left(\ddot\theta_t - \ddot\theta_b\right)\cos\alpha_t
        -
        \left(\dot\theta_t - \dot\theta_b\right)^2\sin\alpha_t
    \right],
\end{equation}
其中 $l_{t,c}$ 为尾电机安装点到尾巴质心的距离。

\subsection{轮地运动学约束}

在尾巴离地工况下，仍假设左右轮满足纯滚动约束和不离地约束，即
\begin{equation}
    a_{wr}^h = R\ddot\theta_{wr}, \qquad
    a_{wl}^h = R\ddot\theta_{wl},
\end{equation}
\begin{equation}
    a_{wr}^v = 0, \qquad a_{wl}^v = 0.
\end{equation}

右腿转轴加速度为
\begin{equation}
    a_r^h =
    R\ddot\theta_{wr}
    + l_r\cos\theta_r \ddot\theta_r
    - l_r\sin\theta_r \dot\theta_r^2,
\end{equation}
\begin{equation}
    a_r^v =
    - l_r\sin\theta_r \ddot\theta_r
    - l_r\cos\theta_r \dot\theta_r^2.
\end{equation}

左腿转轴加速度为
\begin{equation}
    a_l^h =
    R\ddot\theta_{wl}
    + l_l\cos\theta_l \ddot\theta_l
    - l_l\sin\theta_l \dot\theta_l^2,
\end{equation}
\begin{equation}
    a_l^v =
    - l_l\sin\theta_l \ddot\theta_l
    - l_l\cos\theta_l \dot\theta_l^2.
\end{equation}

机体加速度取左右腿转轴加速度平均值：
\begin{equation}
    a_b^h = \frac{a_r^h + a_l^h}{2},
    \qquad
    a_b^v = \frac{a_r^v + a_l^v}{2}.
\end{equation}

偏航角加速度由左右轮水平加速度差得到：
\begin{equation}
    \ddot\phi =
    \frac{a_{wr}^h - a_{wl}^h}{2R_w}
    =
    \frac{R\left(\ddot\theta_{wr} - \ddot\theta_{wl}\right)}{2R_w}.
\end{equation}

定义
\begin{equation}
    \Gamma =
    \frac{
        l_r\cos\theta_r \ddot\theta_r
        +
        l_l\cos\theta_l \ddot\theta_l
    }{2}
    -
    \frac{
        l_r\sin\theta_r \dot\theta_r^2
        +
        l_l\sin\theta_l \dot\theta_l^2
    }{2},
\end{equation}
则有
\begin{equation}
    \ddot X_b^h =
    \frac{R\left(\ddot\theta_{wr}+\ddot\theta_{wl}\right)}{2}
    + \Gamma.
\end{equation}

由 $\ddot X_b^h$ 与 $\ddot\phi$ 可反解左右轮角加速度：
\begin{equation}
    \ddot\theta_{wr}
    =
    \frac{\ddot X_b^h + R_w\ddot\phi}{R}
    -
    \frac{\Gamma}{R},
\end{equation}
\begin{equation}
    \ddot\theta_{wl}
    =
    \frac{\ddot X_b^h - R_w\ddot\phi}{R}
    -
    \frac{\Gamma}{R}.
\end{equation}

通过该变换，轮角加速度不再作为独立广义加速度，而由 $\ddot X_b^h$、$\ddot\phi$、$\ddot\theta_l$ 和 $\ddot\theta_r$ 表示。

\subsection{动力学方程}

基于牛顿--欧拉法，可分别对机体、左右腿、左右轮以及尾巴建立动力学方程。经过消去内部约束力后，可得到六个独立动力学方程：
\begin{equation}
    f_i(q,\dot q,\ddot q,u) = 0,
    \qquad i = 1,2,\cdots,6.
\end{equation}

其中，六个方程分别对应：
\begin{align}
    f_1 &: \text{整体水平方向动量方程}, \\
    f_2 &: \text{机体俯仰转动方程}, \\
    f_3 &: \text{右腿转动方程}, \\
    f_4 &: \text{左腿转动方程}, \\
    f_5 &: \text{偏航转动方程}, \\
    f_6 &: \text{尾巴转动方程}.
\end{align}

下面给出消去内力后的动力学方程显式形式。首先，左右轮的地面水平作用力可由轮子转动方程消去：
\begin{equation}
    F_{g\to wl}^h
    =
    -\frac{T_{wl\to l}+I_{wl}\ddot\theta_{wl}}{R},
    \qquad
    F_{g\to wr}^h
    =
    -\frac{T_{wr\to r}+I_{wr}\ddot\theta_{wr}}{R}.
\end{equation}

尾巴对机体的反作用力由尾巴平动方程给出：
\begin{equation}
    F_{t\to b}^h = -m_t a_t^h,
    \qquad
    F_{t\to b}^v = -m_t g - m_t a_t^v.
\end{equation}

因此，整体水平方向动量方程为
\begin{equation}
\begin{aligned}
    f_1
    ={}&
    m_b a_b^h
    + m_l a_l^h
    + m_r a_r^h
    + m_t a_t^h
    + m_{wl} a_{wl}^h
    + m_{wr} a_{wr}^h  \\
    &+
    \frac{
        T_{wl\to l}
        + T_{wr\to r}
        + I_{wl}\ddot\theta_{wl}
        + I_{wr}\ddot\theta_{wr}
    }{R}
    =0 .
\end{aligned}
\end{equation}

机体俯仰转动方程为
\begin{equation}
\begin{aligned}
    f_2
    ={}&
    I_b\ddot\theta_b
    - T_{l\to b}
    - T_{r\to b}
    - T_{t\to b}  \\
    &-
    \left(
        r_{bp}^h F_{t\to b}^v
        -
        r_{bp}^v F_{t\to b}^h
    \right)
    -
    m_b g l_b \sin(\theta_b+\theta_{b0})
    =0 .
\end{aligned}
\end{equation}

代入尾巴反作用力后，上式也可写为
\begin{equation}
\begin{aligned}
    f_2
    ={}&
    I_b\ddot\theta_b
    - T_{l\to b}
    - T_{r\to b}
    - T_{t\to b} \\
    &-
    \left[
        r_{bp}^h(-m_t g-m_t a_t^v)
        -
        r_{bp}^v(-m_t a_t^h)
    \right]
    -
    m_b g l_b \sin(\theta_b+\theta_{b0})
    =0 .
\end{aligned}
\end{equation}

右腿转动方程为
\begin{equation}
\begin{aligned}
    f_3
    ={}&
    I_r\ddot\theta_r
    - T_{wr\to r}
    + T_{r\to b}
    - m_r g l_{r,d}\sin(\theta_r+\theta_{r0}) \\
    &-
    m_{wr} a_{wr}^h l_r\cos\theta_r
    -
    \frac{
        (T_{wr\to r}+I_{wr}\ddot\theta_{wr})
        l_r\cos\theta_r
    }{R} \\
    &-
    \frac{
        (m_b a_b^v + m_l a_l^v + m_r a_r^v)
        l_r\sin\theta_r
    }{2} \\
    &-
    \frac{
        (m_{wr}-m_b-m_l-m_r-m_{wl})
        g l_r\sin\theta_r
    }{2}
    =0 .
\end{aligned}
\end{equation}

左腿转动方程为
\begin{equation}
\begin{aligned}
    f_4
    ={}&
    I_l\ddot\theta_l
    - T_{wl\to l}
    + T_{l\to b}
    - m_l g l_{l,d}\sin(\theta_l+\theta_{l0}) \\
    &-
    m_{wl} a_{wl}^h l_l\cos\theta_l
    -
    \frac{
        (T_{wl\to l}+I_{wl}\ddot\theta_{wl})
        l_l\cos\theta_l
    }{R} \\
    &-
    \frac{
        (m_b a_b^v + m_l a_l^v + m_r a_r^v)
        l_l\sin\theta_l
    }{2} \\
    &-
    \frac{
        (m_{wl}-m_b-m_l-m_r-m_{wr})
        g l_l\sin\theta_l
    }{2}
    =0 .
\end{aligned}
\end{equation}

偏航转动方程为
\begin{equation}
    f_5
    =
    I_{\mathrm{yaw}}\ddot\phi
    -
    \frac{R_w}{R}
    \left(
        T_{wl\to l}
        + I_{wl}\ddot\theta_{wl}
        - T_{wr\to r}
        - I_{wr}\ddot\theta_{wr}
    \right)
    =0 .
\end{equation}

尾巴转动方程为
\begin{equation}
    f_6
    =
    I_t\ddot\theta_t
    + T_{t \to b}
    - m_t g l_{t,c}\cos\alpha_t
    = 0.
\end{equation}

上述方程仍含有 $\ddot\theta_{wr}$ 和 $\ddot\theta_{wl}$。在后续状态空间建模中，利用轮地运动学约束将其替换为
\begin{equation}
    \ddot\theta_{wr}
    =
    \frac{\ddot X_b^h + R_w\ddot\phi}{R}
    -
    \frac{\Gamma}{R},
    \qquad
    \ddot\theta_{wl}
    =
    \frac{\ddot X_b^h - R_w\ddot\phi}{R}
    -
    \frac{\Gamma}{R}.
\end{equation}
经过该代换后，六个方程仅关于广义坐标 $q$、广义速度 $\dot q$、广义加速度 $\ddot q$ 以及控制输入 $u$。

因此，尾巴通过三种方式影响系统动力学：

\begin{enumerate}
    \item 尾巴质心加速度 $a_t^h$ 参与整体水平方向动量方程；
    \item 尾巴对机体的反作用力通过安装点力臂影响机体俯仰动力学；
    \item 尾电机力矩 $T_{t \to b}$ 直接作用于机体和尾巴转动方程。
\end{enumerate}

\subsection{标准动力学形式}

将六个动力学方程整理为标准矩阵形式：
\begin{equation}
    M(q)\ddot q = B(q)u + g(q,\dot q),
\end{equation}
其中 $M(q)\in\mathbb{R}^{6\times 6}$ 为惯性矩阵，$B(q)\in\mathbb{R}^{6\times 5}$ 为输入矩阵，$g(q,\dot q)$ 包含重力项、离心项、科氏项以及其他非输入项。

在符号计算中，矩阵 $M(q)$ 可由动力学方程对广义加速度求偏导得到：
\begin{equation}
    M_{ij}(q)
    =
    \frac{\partial f_i}{\partial \ddot q_j}.
\end{equation}

输入矩阵由动力学方程对控制输入求偏导得到：
\begin{equation}
    B_{ij}(q)
    =
    -\frac{\partial f_i}{\partial u_j}.
\end{equation}

非输入项可写为
\begin{equation}
    g(q,\dot q)
    =
    -\left[
        f(q,\dot q,\ddot q,u)
        - M(q)\ddot q
        + B(q)u
    \right].
\end{equation}

于是非线性系统可表示为
\begin{equation}
    \ddot q
    =
    M^{-1}(q)
    \left[
        B(q)u + g(q,\dot q)
    \right].
\end{equation}

\subsection{平衡点求解}

为了进行线性化，首先求取系统静态平衡点。设机体和尾巴目标平衡角为
\begin{equation}
    \theta_b^\ast = 0,
    \qquad
    \theta_t^\ast = 0.
\end{equation}

在静态平衡点处，有
\begin{equation}
    \dot q^\ast = 0,
    \qquad
    \ddot q^\ast = 0.
\end{equation}

同时设左右腿对机体的静态力矩相同：
\begin{equation}
    T_{r \to b}^\ast = T_{l \to b}^\ast = T_{\mathrm{leg}}^\ast,
\end{equation}
轮端静态输入为
\begin{equation}
    T_{wr \to r}^\ast = 0,
    \qquad
    T_{wl \to l}^\ast = 0.
\end{equation}

尾巴静态力矩为
\begin{equation}
    T_{t \to b}^\ast = T_t^\ast.
\end{equation}

待求未知量为
\begin{equation}
    \theta_l^\ast,\quad
    \theta_r^\ast,\quad
    T_{\mathrm{leg}}^\ast,\quad
    T_t^\ast.
\end{equation}

选取机体转动方程、左右腿转动方程以及尾巴转动方程作为静态平衡方程：
\begin{equation}
    f_2^\ast = 0,
    \qquad
    f_3^\ast = 0,
    \qquad
    f_4^\ast = 0,
    \qquad
    f_6^\ast = 0.
\end{equation}

联立求解可得到平衡点状态
\begin{equation}
    x^\ast =
    \begin{bmatrix}
        0 &
        0 &
        0 &
        0 &
        \theta_l^\ast &
        0 &
        \theta_r^\ast &
        0 &
        \theta_b^\ast &
        0 &
        \theta_t^\ast &
        0
    \end{bmatrix}^{\mathrm T},
\end{equation}
以及平衡输入
\begin{equation}
    u^\ast =
    \begin{bmatrix}
        T_{\mathrm{leg}}^\ast &
        T_{\mathrm{leg}}^\ast &
        0 &
        0 &
        T_t^\ast
    \end{bmatrix}^{\mathrm T}.
\end{equation}

\subsection{状态空间线性化}

定义一阶状态向量为
\begin{equation}
    x =
    \begin{bmatrix}
        X_b^h &
        \dot X_b^h &
        \phi &
        \dot\phi &
        \theta_l &
        \dot\theta_l &
        \theta_r &
        \dot\theta_r &
        \theta_b &
        \dot\theta_b &
        \theta_t &
        \dot\theta_t
    \end{bmatrix}^{\mathrm T}.
\end{equation}

因此系统状态维数为 $12$，输入维数为 $5$。在平衡点附近定义扰动量
\begin{equation}
    \delta x = x - x^\ast,
    \qquad
    \delta u = u - u^\ast.
\end{equation}

非线性系统在平衡点处一阶线性化后得到
\begin{equation}
    \delta \dot x = A\delta x + B\delta u.
\end{equation}

状态矩阵 $A$ 的运动学部分满足
\begin{equation}
    \dot X_b^h = \dot X_b^h,
    \quad
    \dot\phi = \dot\phi,
    \quad
    \dot\theta_l = \dot\theta_l,
    \quad
    \dot\theta_r = \dot\theta_r,
    \quad
    \dot\theta_b = \dot\theta_b,
    \quad
    \dot\theta_t = \dot\theta_t.
\end{equation}

因此 $A$ 中有如下单位映射：
\begin{equation}
    A_{1,2}=1,\quad
    A_{3,4}=1,\quad
    A_{5,6}=1,\quad
    A_{7,8}=1,\quad
    A_{9,10}=1,\quad
    A_{11,12}=1.
\end{equation}

动力学部分由
\begin{equation}
    \ddot q =
    M^{-1}(q)
    \left[
        B(q)u + g(q,\dot q)
    \right]
\end{equation}
线性化得到。在平衡点处，有
\begin{equation}
    M^\ast = M(q^\ast),
    \qquad
    B^\ast = B(q^\ast).
\end{equation}

记
\begin{equation}
    G_{\theta_i}
    =
    \left.
    \frac{\partial g}{\partial \theta_i}
    \right|_{x=x^\ast,u=u^\ast},
    \qquad
    \theta_i \in
    \left\{
        \theta_l,\theta_r,\theta_b,\theta_t
    \right\}.
\end{equation}

则广义加速度对姿态变量的线性化系数为
\begin{equation}
    A_{\theta_i}^{\mathrm{dyn}}
    =
    \left(M^\ast\right)^{-1}
    G_{\theta_i}.
\end{equation}

输入矩阵的动力学部分为
\begin{equation}
    B^{\mathrm{dyn}}
    =
    \left(M^\ast\right)^{-1}
    B^\ast.
\end{equation}

最终，将 $A_{\theta_i}^{\mathrm{dyn}}$ 填入状态矩阵中对应的加速度行：
\begin{equation}
    \left[
    \begin{array}{c}
        \delta \ddot X_b^h \\
        \delta \ddot \phi \\
        \delta \ddot \theta_l \\
        \delta \ddot \theta_r \\
        \delta \ddot \theta_b \\
        \delta \ddot \theta_t
    \end{array}
    \right]
    =
    \begin{bmatrix}
        A_{\theta_l}^{\mathrm{dyn}} &
        A_{\theta_r}^{\mathrm{dyn}} &
        A_{\theta_b}^{\mathrm{dyn}} &
        A_{\theta_t}^{\mathrm{dyn}}
    \end{bmatrix}
    \begin{bmatrix}
        \delta\theta_l \\
        \delta\theta_r \\
        \delta\theta_b \\
        \delta\theta_t
    \end{bmatrix}
    +
    B^{\mathrm{dyn}}\delta u.
\end{equation}

由此得到完整的线性状态空间模型：
\begin{equation}
    \delta \dot x
    =
    A\delta x + B\delta u.
\end{equation}

\subsection{LQR 控制律}

基于线性化状态空间模型，设计 LQR 控制器。控制律为
\begin{equation}
    \delta u = -K\delta x,
\end{equation}
因此实际控制输入为
\begin{equation}
    u =
    u^\ast - K(x-x^\ast).
\end{equation}

其中 $K$ 为 LQR 增益矩阵，通过求解如下性能指标得到：
\begin{equation}
    J =
    \int_0^\infty
    \left(
        \delta x^{\mathrm T}Q\delta x
        +
        \delta u^{\mathrm T}R\delta u
    \right)
    dt.
\end{equation}

$Q$ 为状态权重矩阵，$R$ 为输入权重矩阵。对于当前尾巴离地模型，$K$ 的维度为
\begin{equation}
    K \in \mathbb{R}^{5\times 12}.
\end{equation}

\end{document}
